\documentclass[aspectratio=169,12pt]{beamer}

%==================================================
% PAQUETES
%==================================================
\usepackage[spanish,es-noshorthands]{babel}
\usepackage[utf8]{inputenc}
\usepackage[T1]{fontenc}
\usepackage{amsmath,amssymb}
\usepackage{xcolor}
\usepackage{tikz}

\usetheme{Madrid}
\setbeamertemplate{navigation symbols}{}

%==================================================
% COLORES AGRONOMÍA
%==================================================
\definecolor{Primary}{HTML}{1B5E20}
\definecolor{Secondary}{HTML}{2E7D32}
\definecolor{Accent}{HTML}{66BB6A}

\setbeamercolor{frametitle}{bg=Primary, fg=white}
\setbeamercolor{title}{fg=Primary}
\setbeamercolor{structure}{fg=Primary}

%==================================================
\title{Expresiones Algebraicas en Agronomía}
\subtitle{Nivel Medio - Avanzado}
\author{MAT016}
\date{}

%==================================================
\begin{document}
	
	%--------------------------------------------------
	\begin{frame}
		\titlepage
	\end{frame}
	
	%--------------------------------------------------
	\begin{frame}{Ruta de la clase}
		\tableofcontents
	\end{frame}
	
	%==================================================
	\section{Contexto Agronómico}
	
	%--------------------------------------------------
	\begin{frame}{Motivación}
		En agronomía, muchos fenómenos se modelan con expresiones algebraicas:
		
		\pause
		
		\begin{itemize}
			\item Producción de cultivos 🌾
			\item Uso de fertilizantes 🌱
			\item Sistemas de riego 💧
		\end{itemize}
		
		\pause
		
		\textbf{Ejemplo:}
		
		\[
		3x^2y + 5x^2y - 2xy^2
		\]
		
		\pause
		
		Representa producción dependiendo de:
		\begin{itemize}
			\item $x$: fertilización
			\item $y$: humedad del suelo
		\end{itemize}
		
	\end{frame}
	
	%==================================================
	\section{Términos y Expresiones}
	
	%--------------------------------------------------
	\begin{frame}{Términos Algebraicos}
		Un término algebraico tiene:
		
		\pause
		
		\begin{itemize}
			\item Coeficiente
			\item Variable
			\item Exponente
		\end{itemize}
		
		\pause
		
		\textbf{Ejemplo agronómico:}
		
		\[
		7x^3y^2
		\]
		
		\pause
		
		Interpretación:
		\begin{itemize}
			\item Fertilización elevada
			\item Alta interacción con humedad
		\end{itemize}
		
	\end{frame}
	
	%--------------------------------------------------
	\begin{frame}{Términos Semejantes}
		\[
		3x^2y + 5x^2y - 2xy^2
		\]
		
		\pause
		
		\begin{itemize}
			\item $3x^2y$ y $5x^2y$ → semejantes
			\item $2xy^2$ → distinto comportamiento agrícola
		\end{itemize}
		
		\pause
		
		\textbf{Resultado:}
		\[
		8x^2y - 2xy^2
		\]
		
	\end{frame}
	
	%==================================================
	\section{Operaciones Algebraicas}
	
	%--------------------------------------------------
	\begin{frame}{Suma y Resta}
		Producción en dos parcelas:
		
		\[
		(4x^2 + 3xy - 2y^2) + (5x^2 - xy + y^2)
		\]
		
		\pause
		
		\[
		= 9x^2 + 2xy - y^2
		\]
		
		\pause
		
		\textbf{Interpretación:}
		\begin{itemize}
			\item Se combinan producciones
			\item Se reducen efectos similares
		\end{itemize}
		
	\end{frame}
	
	%--------------------------------------------------
	\begin{frame}{Multiplicación}
		Aplicación de fertilizante:
		
		\[
		3x(2x^2 - 5xy + y^2)
		\]
		
		\pause
		
		\[
		= 6x^3 - 15x^2y + 3xy^2
		\]
		
		\pause
		
		\textbf{Interpretación:}
		\begin{itemize}
			\item Escalamiento del cultivo
			\item Efectos no lineales
		\end{itemize}
		
	\end{frame}
	
	%==================================================
	\section{Productos Notables}
	
	%--------------------------------------------------
	\begin{frame}{Cuadrado de Binomio}
		\[
		(x + 3y)^2
		\]
		
		\pause
		
		\[
		= x^2 + 6xy + 9y^2
		\]
		
		\pause
		
		\textbf{Interpretación agrícola:}
		\begin{itemize}
			\item $x^2$: crecimiento base
			\item $6xy$: interacción suelo-humedad
			\item $9y^2$: efecto climático fuerte
		\end{itemize}
		
	\end{frame}
	
	%--------------------------------------------------
	\begin{frame}{Diferencia de Cuadrados}
		\[
		(x-4)(x+4)
		\]
		
		\pause
		
		\[
		= x^2 - 16
		\]
		
		\pause
		
		\textbf{Interpretación:}
		\begin{itemize}
			\item Diferencia entre potencial y pérdida
		\end{itemize}
		
	\end{frame}
	
	%==================================================
	\section{Aplicaciones}
	
	%--------------------------------------------------
	\begin{frame}{Ejemplo Aplicado}
		Sistema de riego:
		
		\[
		(2x + y)(x - y)
		\]
		
		\pause
		
		\[
		= 2x^2 - xy - y^2
		\]
		
		\pause
		
		\textbf{Interpretación:}
		\begin{itemize}
			\item Balance entre riego y pérdidas
		\end{itemize}
		
	\end{frame}
	
	%==================================================
	\section{Ejercicios}
	
	%--------------------------------------------------
	\begin{frame}{Ejercicios Guiados}
		Simplifique:
		
		\pause
		
		\[
		(3x + 2y)^2 - (x+y)^2
		\]
		
		\pause
		
		\textbf{Pista:}
		\begin{itemize}
			\item Expanda ambos cuadrados
		\end{itemize}
		
	\end{frame}
	
	%--------------------------------------------------
	\begin{frame}{Ejercicios Propuestos}
		\begin{enumerate}
			\item 
			\[
			(2x^2 + 3xy) + (x^2 - 5xy)
			\]
			
			\pause
			
			\item 
			\[
			(x+2y)(x-2y)
			\]
			
			\pause
			
			\item 
			\[
			(x^2 + xy)(x+y)
			\]
			
			\pause
			
			\item 
			\[
			2x^2 + 3xy + x^2 - xy
			\]
			
		\end{enumerate}
		
	\end{frame}
	
	%==================================================
	\section{Desafío}
	
	%--------------------------------------------------
	\begin{frame}{Desafío Final}
		Modelo de producción:
		
		\[
		P(x,y) = (x+y)^2 + (x-y)^2 + 2xy
		\]
		
		\pause
		
		\textbf{Actividades:}
		\begin{itemize}
			\item Simplifique
			\item Interprete
			\item Detecte redundancias
		\end{itemize}
		
	\end{frame}
	
	%--------------------------------------------------
	\begin{frame}{Cierre}
		\begin{itemize}
			\item Las expresiones modelan fenómenos reales
			\item La simplificación permite entender mejor el sistema
			\item Detectar errores es clave en agronomía
		\end{itemize}
		
	\end{frame}
	
\end{document}